% 出席確認クイズ 第11回　データに基づく意思決定
%   attendance/attendance11.html から生成（解説は本ガイド用に加筆）。

\quizq{1}{「確証バイアス」の説明として適切なものはどれでしょうか?}%
  {損失を過大に嫌う傾向}%
  {\correct{自分の考えを支持する情報ばかりを集めてしまう傾向}}%
  {多数派の意見に合わせる傾向}%
  {最初に見た数値に判断が引きずられる傾向}%
  {正解は (b)。確証バイアスは、自分の仮説に都合のよい情報を重視し、反する情報を軽視してしまう傾向です。(a) は損失回避、(c) は同調（集団への追随）、(d) はアンカリングの説明です。AIに「この考えに反する証拠は」と尋ねるのは、確証バイアスへの対策になります。}

\quizq{2}{「アンカリング」の説明として適切なものはどれでしょうか?}%
  {\correct{最初に提示された数値や情報に判断が引きずられる}}%
  {自分の考えを支持する情報を集める}%
  {直近の出来事を過大評価する}%
  {集団で極端な結論に至る}%
  {正解は (a)。アンカリングは、最初に示された数値や情報（錨）を基準にして、その後の判断が偏ってしまう現象です。(b) は確証バイアス、(c) は利用可能性ヒューリスティック、(d) は集団極性化の説明です。AIが最初に出した数値にも引きずられやすいので注意が必要です。}

\quizq{3}{「Human-in-the-loop」の説明として適切なものはどれでしょうか?}%
  {\correct{AIの判断過程に人が関与し、最終判断や確認を人が行う設計}}%
  {人がまったく関与しない全自動化}%
  {AIを使わないこと}%
  {人がAIの学習データになること}%
  {正解は (a)。Human-in-the-loop は、AIの処理の途中や最終段階に人の確認・判断を組み込む設計の考え方です。(b) は人が関与しない全自動化（human-out-of-the-loop）です。影響の大きい判断ほど人の関与を厚くします。}

\quizq{4}{「データに基づく意思決定(データドリブン)」の説明として適切なものはどれでしょうか?}%
  {データがあれば人の判断は不要}%
  {\correct{勘や経験だけでなく、データの分析結果を根拠に判断する}}%
  {データが多いほど無条件に正しい}%
  {データを無視して直感で決める}%
  {正解は (b)。データドリブンな意思決定は、勘や経験に加えてデータの分析結果を根拠として判断することです。(a)(c) はデータ万能主義で、データの偏りや分析の前提を無視しています。(d) はデータドリブンの逆です。}

\quizq{5}{AIの推奨に従って判断した場合、その結果の責任は誰にあるでしょうか?}%
  {\correct{最終的な判断をした人・組織}}%
  {誰にもない}%
  {AIそのもの}%
  {AIの開発者だけ}%
  {正解は (a)。AIは判断を支援する道具であり、AIの推奨を採用した最終的な責任は判断した人・組織にあります。(c) AI は法的な責任主体ではなく、(d) 開発者の責任は製品の欠陥など限られた場面に限られます。説明責任を果たすために、判断の根拠と経緯を記録しておくことが重要です。}
